\section{Scope and organisation}\label{m:sec:intro}

This chapter assembles the probabilistic and analytic tools used throughout the
thesis.  The emphasis is structural: exponential clocks generate continuous-time
Markov chains; probability generating functions encode branching and compartment
counts; conditioning on non-absorption produces quasi-stationary laws; and the
method of characteristics converts first-order generating-function equations into
explicit solutions.  Population, absorption and catastrophe are treated as
abstract stochastic mechanisms rather than tied to a particular application.

Two recurring distinctions organise the discussion.  The first is between
ordinary moments and quantities conditioned on survival.  In a subcritical process
the unconditional mean tends to zero, while the population seen conditional on
survival approaches a non-degenerate scale.  The second is between discrete and
continuous time.  Both descriptions may share an embedded event chain, yet the
availability of an explicit time parameter can make their survival constants
analytically quite different.

\Cref{m:sec:prelim} develops exponential clocks, competing events, generators,
Kolmogorov equations and generating functions, ending with the linear
birth--death process.  \Cref{m:sec:gw} then treats the binary Galton--Watson
process, including extinction, the critical lifetime and total progeny.
\Cref{m:sec:qsd} develops conditioning on late survival and introduces the discrete
constant
\[
  A(p)=\lim_{n\to\infty}\frac{S_n}{(2p)^n}
\]
only to the extent required for the limiting conditional moments.

\Cref{m:sec:small} turns to small founding populations, catastrophe exposure and
the continuous-time birth--death comparison.  There the closed form
\(A_{\mathrm c}(p)=(1-2p)/(1-p)\) is derived before it is contrasted with the
discrete constant.  The product, series, bounds, near-critical analysis,
closed-form search and Koenigs-function theory of \(A(p)\) are deliberately reserved
for \ChA.  Finally, \cref{m:sec:moc} develops a ladder of absorption models and
solves their bivariate generating-function equations by characteristics; the
appendices record the supporting hypergeometric identity and coefficient
extraction.

This is the Pass-1 content architecture.  It retains the full harvested toolkit
spine; a later structural pass may rearrange or compress it, but no such redesign is
undertaken here.
